\chapter{Municipality Summary}
\label{app:summary}

\section{Problem or decision addressed}

\noindent\fbox{%
    \parbox{\dimexpr\textwidth-2\fboxsep-2\fboxrule\relax}{%
        Utrecht has to decide where to inspect and intervene first across more
        than 60\,km of quay walls, and which measure to take (monitor, restrict
        use, strengthen or replace) until a wall is renewed, within a limited
        budget and capacity and without bringing the city to a standstill. It
        wants to stop programming on the theoretical end of life. Our method
        supports that decision by ranking walls on risk, the probability of
        failure and its consequences, with the uncertainty of the evidence made
        explicit.
    }%
}

\section{Output: type of decision support}

\noindent\fbox{%
    \parbox{\dimexpr\textwidth-2\fboxsep-2\fboxrule\relax}{%
        A traceable screening and ranking method with its code and data. For
        each wall it gives a risk class (red, orange or green) from a relative
        probability class and a consequence class, the measure that fits the
        class, and the measurement that would most change its class. For the
        test case it delivers a map and a ranked list of the 17 walmuren still on
        timber piles (Figure~\ref{fig:results-map}), the expected effect and cost
        of seven interventions (Section~\ref{sec:tradeoff}) and the risk carried
        until renewal at different paces of the programme
        (Section~\ref{sec:regression}). The level of detail is the single wall
        or rak.
    }%
}

\section{Area analysed}

\noindent\fbox{%
    \parbox{\dimexpr\textwidth-2\fboxsep-2\fboxrule\relax}{%
        The method covers the whole portfolio. It is tested on the walmuren of
        the Wervengebied along the Oudegracht, Nieuwegracht and Kromme
        Nieuwegracht: 51 walls, 4{,}062\,m in 27 raks, of which 17 walls (611\,m)
        still stand on their original timber piles with a werf. This area is used
        because only here drawings, dive records and renewal data exist to build
        and check a failure model; the wharf cellars, the street and the bridges
        are outside the system but set the consequences.
    }%
}

\section{Approach}

\noindent\fbox{%
    \parbox{\dimexpr\textwidth-2\fboxsep-2\fboxrule\relax}{%
        (1) Check the asset register: wrong, collective or rounded years,
        widths with different meanings, and renewals that are not recorded.
        (2) Renewed or recent walls are green and stay in regular inspection;
        the end of life is not used as a gate. (3) For masonry walls on timber
        piles with a werf, compute the failure probability by Monte Carlo
        simulation from the measured cross-sections, updated with the fact that
        the wall stands (proven strength); for other walls, score indicators:
        satellite movement, NEN~2767 condition, maintenance, age and data
        quality. (4) Combine the relative probability class with a consequence
        class (people on the werf, cellars, reachability) into red, orange or
        green. (5) Compare measures in a fault tree and a decision tree, and
        extrapolate the pace of renewal by regression.
    }%
}

\section{Key limitations and assumptions}

\noindent\fbox{%
    \parbox{\dimexpr\textwidth-2\fboxsep-2\fboxrule\relax}{%
        The failure model rests on two schematic cross-sections and literature
        values for the timber and the soil. It is conservative for the older wall
        type and leaves out settlement from pile rot and loss of the timber
        screen, so its results are a ranking and an order of magnitude, not a
        safety verdict under NEN~8700; that needs calibration with the wall type
        and pile condition of a sample of walls. Most human and organisational
        probabilities and the costs are team estimates, as no municipal records
        were used. The consequence score leaves out cables, pipes and trees, and
        the indicator score for walls without a model has not yet been applied
        to the rest of the portfolio.
    }%
}

\section{Why should the Municipality of Utrecht select this approach}

\noindent\fbox{%
    \parbox{\dimexpr\textwidth-2\fboxsep-2\fboxrule\relax}{%
        You want to move from a theoretical end of life to the real condition of
        your walls, but the evidence is incomplete. This approach uses what you
        already hold, the register, renewal records and open data, and shows for
        each wall how uncertain its class is and which measurement would settle
        it. In the test case it found that 11 of the 29 walmuren shown as past
        their design life were already renewed, that the unknown wall type is the
        evidence worth buying first, and that for walls due for renewal within a
        few years, keeping the edge of the werf free is the cheapest interim
        measure, while monitoring pays for walls that will wait longer. Every
        step is traceable, so choices can be explained to the council. It does
        not replace the structural assessment of a single wall, but it gives a
        defensible order in which to spend limited budget and capacity.
    }%
}
